\documentclass[11pt,a4paper]{article}

\usepackage[utf8]{inputenc}
\usepackage[T1]{fontenc}
\usepackage{lmodern}

\usepackage{amsmath}
\usepackage{amssymb}
\usepackage{amsthm}
\usepackage{bm}

\usepackage{booktabs}
\usepackage{geometry}

\geometry{margin=2.5cm}

\title{Loss Functions for Stellar Spectrum Reconstruction}
\author{Yaiza Mar\'ia Mart\'in Mesones}
\date{\today}

\begin{document}

\maketitle

\section{Loss Functions}

The choice of the loss function plays a fundamental role in the reconstruction of high-resolution stellar spectra from low-resolution Gaia observations. Let

\[
x\in\mathbb{R}^{N}
\]

denote the input Gaia spectrum,

\[
y\in\mathbb{R}^{N}
\]

the corresponding reference CALSPEC spectrum, and

\[
\hat{y}=f_{\theta}(x)
\]

the prediction produced by a neural network with parameters
$\theta$.

Here,

\begin{itemize}
    \item $N$ is the number of wavelength samples,
    \item $y_i$ denotes the reference flux at wavelength $\lambda_i$,
    \item $\hat{y}_i$ denotes the predicted flux at wavelength $\lambda_i$,
    \item $\Delta\lambda_i$ represents the wavelength spacing.
\end{itemize}

%%%%%%%%%%%%%%%%%%%%%%%%%%%%%%%%%%%%%%%%%%%%%%%%%%%%%%%%%%%%

\subsection{Mean Squared Error (MSE)}

The Mean Squared Error is defined as

\begin{equation}
L_{\mathrm{MSE}}
=
\frac{1}{N}
\sum_{i=1}^{N}
\left(y_i-\hat{y}_i\right)^2.
\end{equation}

This objective strongly penalizes large reconstruction errors and is one of the standard loss functions for regression problems.

%%%%%%%%%%%%%%%%%%%%%%%%%%%%%%%%%%%%%%%%%%%%%%%%%%%%%%%%%%%%

\subsection{Mean Absolute Error (MAE)}

The Mean Absolute Error is defined as

\begin{equation}
L_{\mathrm{MAE}}
=
\frac{1}{N}
\sum_{i=1}^{N}
|y_i-\hat{y}_i|.
\end{equation}

Unlike MSE, MAE provides a linear penalty and is therefore more robust to outliers.

%%%%%%%%%%%%%%%%%%%%%%%%%%%%%%%%%%%%%%%%%%%%%%%%%%%%%%%%%%%%

\subsection{Huber Loss}

The Huber loss combines the advantages of MSE and MAE.

Define the reconstruction error

\begin{equation}
e_i=y_i-\hat{y}_i.
\end{equation}

Then,

\begin{equation}
L_{\mathrm{Huber}}
=
\frac{1}{N}
\sum_{i=1}^{N}
\ell_{\delta}(e_i),
\end{equation}

where

\begin{equation}
\ell_{\delta}(e)=
\begin{cases}
\dfrac{1}{2}e^{2},
&
|e|<\delta,
\\[0.3cm]
\delta\left(|e|-\dfrac{\delta}{2}\right),
&
|e|\geq\delta,
\end{cases}
\end{equation}

and $\delta>0$ is a hyperparameter controlling the transition between quadratic and linear behavior.

%%%%%%%%%%%%%%%%%%%%%%%%%%%%%%%%%%%%%%%%%%%%%%%%%%%%%%%%%%%%

\section{Physics-Informed Loss Function}

Since stellar spectra satisfy physical constraints, additional regularization terms can be incorporated into the optimization objective.

We propose the following composite loss:

\begin{equation}
L
=
L_{\mathrm{Huber}}
+
\alpha L_{\mathrm{flux}}
+
\beta L_{\mathrm{grad}},
\end{equation}

where

\begin{itemize}
    \item $\alpha\ge0$ controls the importance of flux conservation,
    \item $\beta\ge0$ controls the importance of gradient preservation.
\end{itemize}

%%%%%%%%%%%%%%%%%%%%%%%%%%%%%%%%%%%%%%%%%%%%%%%%%%%%%%%%%%%%

\subsection{Flux Conservation Loss}

The integrated flux of a spectrum is

\begin{equation}
F
=
\int_{\lambda_{\min}}^{\lambda_{\max}}
y(\lambda)\,d\lambda.
\end{equation}

For discretized spectra,

\begin{equation}
F
\approx
\sum_{i=1}^{N}
y_i\Delta\lambda_i.
\end{equation}

The corresponding loss term is therefore

\begin{equation}
L_{\mathrm{flux}}
=
\left(
\sum_{i=1}^{N}
y_i\Delta\lambda_i
-
\sum_{i=1}^{N}
\hat{y}_i\Delta\lambda_i
\right)^2.
\end{equation}

This term encourages the reconstructed spectrum to preserve the total emitted energy.

%%%%%%%%%%%%%%%%%%%%%%%%%%%%%%%%%%%%%%%%%%%%%%%%%%%%%%%%%%%%

\subsection{Gradient Preservation Loss}

Absorption and emission lines are characterized by local variations of the spectrum.

Approximating the derivative by finite differences,

\begin{equation}
\frac{dy}{d\lambda}
\approx
\frac{y_{i+1}-y_i}{\Delta\lambda_i},
\end{equation}

the gradient loss is defined as

\begin{equation}
L_{\mathrm{grad}}
=
\frac{1}{N-1}
\sum_{i=1}^{N-1}
\left[
\left(
\frac{y_{i+1}-y_i}{\Delta\lambda_i}
\right)
-
\left(
\frac{\hat{y}_{i+1}-\hat{y}_i}{\Delta\lambda_i}
\right)
\right]^2.
\end{equation}

This regularization penalizes discrepancies in the local spectral slope and encourages the preservation of line profiles.

%%%%%%%%%%%%%%%%%%%%%%%%%%%%%%%%%%%%%%%%%%%%%%%%%%%%%%%%%%%%

\section{Experimental Comparison}

The proposed objective can be evaluated against standard regression losses using identical neural network architectures.

\begin{table}[h]
\centering
\begin{tabular}{lcccc}
\toprule
Loss Function &
RMSE &
MAE &
Flux Conservation &
Line Reconstruction \\
\midrule
MSE & & & & \\
MAE & & & & \\
Huber & & & & \\
Proposed Loss & & & & \\
\bottomrule
\end{tabular}
\caption{Comparison of different loss functions for stellar spectrum reconstruction.}
\end{table}



\section{Final Objective}

The complete optimization problem can be written as

\begin{equation}
\boxed{
\theta^{*}
=
\arg\min_{\theta}
\left(
L_{\mathrm{Huber}}
+
\alpha L_{\mathrm{flux}}
+
\beta L_{\mathrm{grad}}
\right),
}
\end{equation}

where $\theta^{*}$ denotes the optimal neural network parameters.

This formulation combines a robust regression loss with physically motivated regularization terms designed to preserve both the integrated flux and the local spectral morphology.

\end{document}